\section{Parameters and Integration}
\label{sec:parameters}

This section explains how parameters interact across the ARC-SPRUCE stack: the target-shaped commitment, the
BB-tran rational hash, the lattice preimage sampler, SmallWood, and the PRF.  Detailed parameter discussion and
benchmark context are deferred to Appendix~\ref{app:ext-params}.

\subsection{Base ring and primary profile}
All components share the same base ring \(\Rq\) and modulus \(q\): the commitment matrices, signature map
\(\mA\), and hash key \(B\) are defined over \(\Rq\); SmallWood constraints are evaluated over \(\Fq\); and the PRF
is instantiated over \(\Fq\) so that its constraints are native to the proof system.  The primary parameter profile is
\[
  N=512,
  \qquad
  q=1054721,
  \qquad
  \ell_M=1,
  \qquad
  k_s=2,
  \qquad
  n_c=1,
  \qquad
  \BoundMsg=8.
\]
The condition \(\ell_M=1\) is an encoding assumption for \(M=\vecm\Vert\veck\).  If the semantic message does not
fit into one ring polynomial, \(\ell_M\) must be increased and the commitment dimensions, binding length, MLWE/MSIS
estimates, and SmallWood row inventories must be recomputed.

\subsection{Commitment parameters}
The commitment term is
\[
  c=C_MM+A_s s+e\in\Rq^{n_c},
\]
where
\[
  C_M\in\Rq^{n_c\times \ell_M},
  \qquad
  A_s\in\Rq^{n_c\times k_s},
  \qquad
  s\in\R^{k_s},
  \qquad
  e\in\R^{n_c}.
\]
For the primary profile,
\[
  C_M\in\Rq^{1\times1},
  \qquad
  A_s\in\Rq^{1\times2},
  \qquad
  s\in\R^2,
  \qquad
  e\in\R.
\]
Thus the commitment opening contains three small ring polynomials beyond the semantic message: two secret
polynomials in \(s\) and one error polynomial \(e\).  The coefficients of \(s\) and \(e\) are sampled uniformly from
\([-8,8]\).  For estimator reporting, this distribution has coefficient standard deviation
\[
  \sigma_B=\sqrt{B(B+1)/3}=\sqrt{24}
  \qquad\text{for }B=8.
\]
The proof bound \(B=8\), the estimator standard deviation \(\sigma_B\), the signature preimage bound \(\BoundSig\),
and the message/key bound \(\BoundMsg\) are distinct quantities; the first parameter profile sets \(B=\BoundMsg=8\)
for the opened commitment coordinates.

The binding matrix is
\[
  [C_M\mid A_s\mid I_{n_c}],
\]
so the binding input length is
\[
  L_{\mathsf{bind}}=\ell_M+k_s+n_c=1+2+1=4.
\]
For double openings with all opened coefficients bounded by \(\BoundMsg=8\), the coefficient-infinity bound on the
difference vector is
\[
  \beta_{\mathsf{bind}}=2\BoundMsg=16.
\]
The hiding claim is computational under Assumption~\ref{ass:mlwe-hiding}; there is no information-theoretic
commitment-hiding claim.

\subsection{BB-tran rational-hash parameters}
The DKLW BB-tran part uses issuer-sampled values
\[
  \ell_{\mu_{\mathsf{sig}}}=1,
  \qquad
  \ell_{x_0}=2,
  \qquad
  x_1\in\Rq.
\]
The target contribution is
\[
  h_{\muSig,\chiSig}(B)=B_0+B_1\muSig+B_2x_0+Z,
  \qquad
  (B_3-\mathbf 1_nx_1)\Had Z=1_n.
\]
The complete signed target is
\[
  T=B_0+B_1\muSig+B_2x_0+Z+C_MM+A_s s+e.
\]
Honest target derivation restarts when \(B_3-\mathbf 1_nx_1\notin\Rq^{\times n}\).  The restart probability must be
negligible or explicitly carried as the rational-hash admissibility error \(\delta\).  The semantic message \(M\)
never appears in the \(B_1\) rational-hash term.

\paragraph{Non-blind reduction parameters.}
The non-blind reduction imported from DKLW is parameterized by a norm bound \(\beta_{\mathrm{DKLW}}\), an
inadmissibility bound \(\delta\), and the predicate/linear-independence side conditions of the BB-tran family.  The
public verifier bound \(\BoundSig\) should dominate the relevant DKLW norm bound under the chosen norm
comparison.  The IntGenISIS lifting is cited only under its stated side conditions and is not a complete proof of the
interactive ARC issuance layer.

\subsection{Sampler parameters}
The trapdoor sampler must output short signatures \(\vecu\) with overwhelming probability and with a distribution
close to the required lattice-coset distribution for the imported signature theorem.  This requires choosing the ring
degree and modulus, trapdoor quality, Gaussian widths, and an acceptance predicate that yields the public
verification bound \(\BoundSig\).  Appendix~\ref{app:gaussian-sampler} describes the NTRU/Antrag-based candidate
used for tuning.  The commitment design does not change the sampler/trapdoor admissibility interface.

\subsection{SmallWood row arithmetic}
The pre-sign proof checks the commitment opening, message encoding, key-space membership, bounds, and policy
constraints.  The showing proof checks the signature equation, the bilinear inverse equation, the hidden commitment
term, the message encoding, and the PRF checkpoint constraints.

Ignoring PRF internals and auxiliary range rows, the primary profile has the following showing witness inventory:
\[
  \vecu:2,
  \quad
  M:1,
  \quad
  s:2,
  \quad
  e:1,
  \quad
  \muSig:1,
  \quad
  x_0:2,
  \quad
  x_1:1,
  \quad
  Z:1.
\]
This is a total of
\[
  2+1+2+1+1+2+1+1=11
\]
ring polynomials.  With SmallWood packing \(s_{\mathrm{SW}}=16\), one ring polynomial contributes
\[
  N/s_{\mathrm{SW}}=512/16=32
\]
witness rows.  Therefore the non-PRF showing witness inventory contributes
\[
  11\cdot 32=352
\]
SmallWood rows before PRF auxiliary rows.  The pre-sign opening inventory is \((M,s,e)\), namely
\[
  1+2+1=4
\]
ring polynomials, contributing \(4\cdot32=128\) rows at the same packing.  These are row-count arithmetic
statements only; final proof sizes require running the SmallWood optimizer with the actual row inventory, constraint
degrees, Fiat--Shamir parameters, and PRF checkpoint schedule.

\subsection{Secondary compact profile}
A smaller-ring candidate is
\[
  N=256,
  \qquad
  q=1054721,
  \qquad
  \ell_M=1,
  \qquad
  k_s=4,
  \qquad
  n_c=1,
  \qquad
  B=8.
\]
Its MLWE hiding estimate is approximately \(194.4\) bits for the estimator configuration used in the parameter
notes.  The corresponding estimator dimensions are
\[
  m_{\mathsf{LWE}}=Nn_c=256,
  \qquad
  m_{\mathsf{SIS}}=N(\ell_M+k_s+n_c)=256(1+4+1)=1536.
\]
Its non-PRF showing witness inventory is
\[
  \vecu:2,
  \quad M:1,
  \quad s:4,
  \quad e:1,
  \quad \muSig:1,
  \quad x_0:2,
  \quad x_1:1,
  \quad Z:1,
\]
for a total of \(13\) ring polynomials.  With \(s_{\mathrm{SW}}=16\), this contributes
\[
  13\cdot(256/16)=208
\]
SmallWood rows before PRF auxiliary rows.  A binding estimate must be inserted before this profile is claimed as a
full concrete commitment-security profile.

\subsection{PRF parameters}
The PRF must be secure at the target security level, efficiently representable in constraints, and compatible with
\(\Fq\).  We use a Poseidon2-style permutation with feed-forward and truncation.  Its parameters include the state
width \(t\), round counts \((R_F,R_P)\), MDS matrices, round constants, S-box exponent \(d\), truncation length
\(\ell_{\mathsf{tag}}\), and key space \(\mathcal K_{\mathsf{key}}\).  We require
\(\ell_{\mathsf{tag}}\ge\lceil\lambda/\log_2(q)\rceil\) as an output-length condition; pseudorandomness of the
resulting tags is exactly Assumption~\ref{ass:poseidon-prf} for keys sampled from
\(\mathcal K_{\mathsf{key}}\).  The output-length condition alone is not a proof of PRF security.
